% Generated draft: requires mathematical and source review.
\documentclass[11pt]{article}
\usepackage[margin=25mm]{geometry}
\usepackage{fontspec}
\setmainfont{Latin Modern Roman}
\usepackage{amsmath,amssymb,mathtools,unicode-math}
\setmathfont{Latin Modern Math}
\usepackage{enumitem,needspace,xurl,hyperref}
\setlength{\parindent}{0pt}
\setlength{\parskip}{5pt}
\setlength{\emergencystretch}{3em}
\newcommand{\N}{\mathbb{N}}
\newcommand{\Z}{\mathbb{Z}}
\newcommand{\Q}{\mathbb{Q}}
\newcommand{\R}{\mathbb{R}}
\newcommand{\C}{\mathbb{C}}
\newcommand{\F}{\mathbb{F}}
\newcommand{\A}{\mathbb{A}}
\newcommand{\PP}{\mathbb{P}}
\newcommand{\E}{\mathbb{E}}
\newcommand{\eps}{\varepsilon}
\newcommand{\dd}{\,\mathrm{d}}
\begin{document}
\section*{UnsolvedMath ID: 700144}
\textbf{Source problem code:} catalogue entry
\subsection{Definitions and mathematical statement}
Let \(G\) be a finite group with identity \(1\). An involution is an element \(t\in G\) of order exactly \(2\), equivalently \(t\ne 1\) and \(t^2=1\). Its centralizer is
\[
C_G(t)=\{g\in G:gt=tg\}.
\]
For a group \(H\), define its derived series by \(H^{(0)}=H\) and
\[
H^{(n+1)}=[H^{(n)},H^{(n)}]\qquad(n\geq 0),
\]
where \([K,K]\) is the subgroup generated by all \(x^{-1}y^{-1}xy\) with \(x,y\in K\). The group \(H\) is soluble if \(H^{(m)}=\{1\}\) for some integer \(m\geq 0\).

Describe, up to isomorphism, all finite groups \(G\) satisfying
\[
\forall t\in G,\qquad
\bigl(t\ne 1\ \text{and}\ t^2=1\bigr)
\Longrightarrow
\bigl(\exists m\in\Z_{\geq 0}: C_G(t)^{(m)}=\{1\}\bigr).
\]
The integer \(m\) may depend on \(t\). The source imposes no additional hypotheses on \(G\); in particular, the condition holds vacuously when \(G\) has no involutions. The source calls this a ``well-known problem.'' Its supplied open status is a registry label, not an independently verified assessment.
\subsection{Short English statement}
Classify the finite groups in which the centralizer of every element of order exactly \(2\) is soluble. Groups without such elements are included.
\subsection{Sources}
\begin{itemize}
\item Supplied local UnsolvedMath registry record: permanent ID \(700144\), exact source problem code \texttt{catalogue entry}, title \emph{Kourovka Notebook Problem 2.74}; upstream revision \texttt{d371ee917ac9647fbb06cd6fc529deec6f3dda82}.
\item The record attributes the problem to A. I. Starostin, proposed in 1967, and identifies its source as \emph{The Kourovka Notebook (Unsolved Problems in Group Theory)}, second issue (1967), problem 2.74, as included in the twenty-first issue edition (2026), \texttt{arXiv:1401.0300v48}.
\item The supplied collection description names E. I. Khukhro and V. D. Mazurov as editors. These bibliographic details and the registry's open label are supplied source information and have not been independently verified.
\end{itemize}
\end{document}
